% ===================================================================
% commands.tex — تنظیمات و دستورات کمکی
% ===================================================================

% --- بسته‌های ریاضی ---
\usepackage{amssymb,amsmath,amsfonts}

% --- حاشیه‌های صفحه ---
\usepackage[top=30mm, bottom=30mm, left=25mm, right=30mm]{geometry}

% --- تصاویر و رنگ ---
\usepackage{graphicx}
\usepackage{color}
\usepackage{xcolor}

% --- فاصله‌ی عمودی خطوط ---
\usepackage{setspace}

% --- لینک‌ها ---
\usepackage[pagebackref=false,hidelinks]{hyperref}

% --- سرصفحه ---
\usepackage{fancyhdr}

% --- سایر ---
\usepackage{verbatim}

% --- جداول ---
\usepackage{booktabs}
\usepackage{longtable}
\usepackage{multirow}
\usepackage{tabularx}
\usepackage{array}

% --- شناور ---
\usepackage{float}
\usepackage{caption}
\usepackage{subcaption}

% --- قطعه‌ی کد ---
\usepackage{listings}

% ===================================================================
% فراخوانی xepersian — حتماً آخرین بسته باشد
% ===================================================================
\usepackage{xepersian}
\SepMark{-}
\let\cleardoublepage\clearpage

% --- فونت‌های فارسی ---
\settextfont[Scale=1.2]{B Nazanin}
\setlatintextfont{Times New Roman}
\renewcommand{\labelitemi}{$\bullet$}
\setlength{\headheight}{15pt}
\emergencystretch=2em

% --- فونت اعداد در فرمول ---
\setdigitfont[Scale=1.1]{Times New Roman}

% --- فونت‌های اضافی ---
\defpersianfont\nastaliq[Scale=2]{B Nazanin}
\defpersianfont\chapternumber[Scale=3]{B Nazanin}

% ===================================================================
% دستورات کمکی اختصاصی این پروژه
% ===================================================================

% اصطلاح فارسی با پاورقی انگلیسی:  \term{فارسی}{English}
\newcommand{\term}[2]{#1\RTLfootnote{\lr{#2}}}

% کد یک‌خطی:  \code{text}
\newcommand{\code}[1]{\texttt{\lr{\detokenize{#1}}}}

% ===================================================================
% تنظیمات قطعه‌ی کد
% ===================================================================
\definecolor{codegray}{rgb}{0.5,0.5,0.5}
\definecolor{codepurple}{rgb}{0.58,0,0.82}
\definecolor{backcolour}{rgb}{0.97,0.97,0.97}
\definecolor{codeblue}{rgb}{0.0,0.3,0.7}

\lstdefinestyle{phpstyle}{
  backgroundcolor   = \color{backcolour},
  commentstyle      = \color{codegray}\itshape,
  keywordstyle      = \color{codeblue}\bfseries,
  numberstyle       = \tiny\color{codegray},
  stringstyle       = \color{codepurple},
  basicstyle        = \ttfamily\footnotesize,
  breakatwhitespace = false,
  breaklines        = true,
  captionpos        = b,
  keepspaces        = true,
  numbers           = left,
  numbersep         = 5pt,
  showspaces        = false,
  showstringspaces  = false,
  showtabs          = false,
  tabsize           = 2,
  frame             = single,
  rulecolor         = \color{gray!40},
  language          = PHP
}

\lstdefinestyle{jsonStyle}{
  backgroundcolor   = \color{backcolour},
  basicstyle        = \ttfamily\footnotesize,
  breaklines        = true,
  captionpos        = b,
  numbers           = left,
  numberstyle       = \tiny\color{codegray},
  numbersep         = 5pt,
  frame             = single,
  rulecolor         = \color{gray!40},
}

\lstset{style=phpstyle}
\renewcommand{\lstlistingname}{کد}
\renewcommand{\lstlistlistingname}{فهرست کدها}
\makeatletter
\@addtoreset{lstlisting}{chapter}
\makeatother
\renewcommand{\thelstlisting}{\thechapter-\arabic{lstlisting}}
\newcommand{\listingcaption}[1]{%
  \refstepcounter{lstlisting}%
  \par\vspace{4pt}%
  \begin{center}
    \small \textbf{\lstlistingname~\thelstlisting:} #1
  \end{center}
  \vspace{6pt}%
}

% ===================================================================
% تنظیمات تکمیلی
% ===================================================================
\renewcommand{\bibname}{منابع و مراجع}

\newcommand\tocheading{\par عنوان\hfill صفحه \par}
\newcommand\lofheading{\hspace*{.5cm}\figurename\hfill صفحه \par}
\newcommand\lotheading{\hspace*{.5cm}\tablename\hfill صفحه \par}

% --- دستورهای سفارشی‌سازی صفحات اول فصل‌ها ---
\makeatletter
\newcommand\mycustomraggedright{%
 \if@RTL\raggedleft%
 \else\raggedright%
 \fi}
\def\@makechapterhead#1{%
\thispagestyle{style1}
\vspace*{20\p@}%
{\parindent \z@ \mycustomraggedright
\ifnum \c@secnumdepth >\m@ne
\if@mainmatter
\bfseries{\Huge \@chapapp}\small\space {\chapternumber\thechapter}
\par\nobreak
\vskip 0\p@
\fi
\fi
\interlinepenalty\@M 
\Huge \bfseries #1\par\nobreak
\vskip 40\p@
}
}
\bidi@patchcmd{\@makechapterhead}{\thechapter}{\tartibi{chapter}}{}{}
\bidi@patchcmd{\chaptermark}{\thechapter}{\tartibi{chapter}}{}{}
\makeatother

% --- سبک‌های سرصفحه ---
\pagestyle{fancy}
\renewcommand{\chaptermark}[1]{\markboth{\chaptername~\tartibi{chapter}: #1}{}}

\fancypagestyle{style1}{
\fancyhf{} 
\fancyfoot[c]{\thepage}
\fancyhead[R]{\leftmark}
\renewcommand{\headrulewidth}{1.2pt}
}

\fancypagestyle{style2}{
\fancyhf{}
\fancyhead[R]{چکیده}
\fancyfoot[C]{\thepage{}}
\renewcommand{\headrulewidth}{1.2pt}
}

\fancypagestyle{style3}{
  \fancyhf{}
  \fancyhead[R]{فهرست نمادها}
  \fancyfoot[C]{\thepage}
  \renewcommand{\headrulewidth}{1.2pt}
}

\fancypagestyle{style4}{
  \fancyhf{}
  \fancyhead[R]{فهرست جداول}
  \fancyfoot[C]{\thepage}
  \renewcommand{\headrulewidth}{1.2pt}
}

\fancypagestyle{style5}{
  \fancyhf{}
  \fancyhead[R]{فهرست اشکال}
  \fancyfoot[C]{\thepage}
  \renewcommand{\headrulewidth}{1.2pt}
}

\fancypagestyle{style6}{
  \fancyhf{}
  \fancyhead[R]{فهرست مطالب}
  \fancyfoot[C]{\thepage}
  \renewcommand{\headrulewidth}{1.2pt}
}

\fancypagestyle{style7}{
  \fancyhf{}
  \fancyhead[R]{نمایه}
  \fancyfoot[C]{\thepage}
  \renewcommand{\headrulewidth}{1.2pt}
}

\fancypagestyle{style8}{
  \fancyhf{}
  \fancyhead[R]{منابع و مراجع}
  \fancyfoot[C]{\thepage}
  \renewcommand{\headrulewidth}{1.2pt}
}

\fancypagestyle{style9}{
  \fancyhf{}
  \fancyhead[R]{واژه‌نامه‌ی فارسی به انگلیسی}
  \fancyfoot[C]{\thepage}
  \renewcommand{\headrulewidth}{1.2pt}
}

% --- دستور حذف نام لیست از فهرست مطالب ---
\newcommand*{\BeginNoToc}{%
  \addtocontents{toc}{%
    \edef\protect\SavedTocDepth{\protect\the\protect\value{tocdepth}}%
  }%
  \addtocontents{toc}{%
    \protect\setcounter{tocdepth}{-10}%
  }%
}
\newcommand*{\EndNoToc}{%
  \addtocontents{toc}{%
    \protect\setcounter{tocdepth}{\protect\SavedTocDepth}%
  }%
}
\newcounter{savepage}
\renewcommand{\listfigurename}{فهرست اشکال}
\renewcommand{\listtablename}{فهرست جداول}
